\section{Similarity search and alignment results}
\label{sec:appendix-search}

\texttt{mmsearch(target, query)} requests sequence search from a custom scan.
The first operand is the stored target or a supported transformation of it;
the second is the query. This section extends the quick-start example with
multiple matches, result fields, thresholds, prepared statements, and
controlled comparisons of execution routes.

\subsection{Prepare a small target table}

The following targets combine ordinary metadata with DNA. Target 5 contains
three exact copies of the query used below. The examples deliberately keep
alignment rows separate from table rows.

\par\noindent\begin{minipage}{\linewidth}
\begin{lstlisting}[style=bioseq-appendix-sql]
SET bioseq.index.insert_maintenance.enabled = on;
CREATE TABLE tutorial_mmsearch (
    pk integer PRIMARY KEY,
    organism text NOT NULL,
    collected_on date NOT NULL,
    seq0 bioseq.seq('dna4', 64)
) USING bioseq;
INSERT INTO tutorial_mmsearch VALUES
    (1, 'Homo sapiens', DATE '2026-06-10',
     'TCACTGTAGGATCGTAGCTAGTCCGTACTGTAGCTGATCGTAGC'
     'CCTAGTCACTGTAGCCTGATCGTAGCTAGTCACTGTAGCCATACCTGAA'),
    (2, 'Homo sapiens', DATE '2025-11-04',
     'CCCCAACTTAGTCACTTCTCAAAAATTACTCTCTCATCATCATC'
     'TACTCCTAGTCACTGTAGCCTGATCGTAGCTAGTCACACACACCCCAAA'),
    (3, 'Homo sapiens', DATE '2024-02-18',
     'GGGGGTCACTTCTCAAAAATTAGGGGGGGGGTAGTCACTGTAG'
     'CCTGATCGTAGCTAGGACACATATCATCATCTACTCCTATCATCATCTACTCCTA'),
    (4, 'Mus musculus', DATE '2026-07-01',
     'TTTCATCATCATCATCTACTCCTATCTATCATCATCATCATCTACTCCTA'
     'TCTACTCTCATCATCCTGATCGTAGCTATACTCCTACTATCCTATT'),
    (5, 'Homo sapiens', DATE '2026-01-20',
     'ACGTCTAGTCACTGTAGCCTGATCGTAGCTAGTGTACTCTAGTCACTGTA'
     'GCCTGATCGTAGCTAGTTGATCCCCATATCTTCACTAGTCACTGTAGCCT'
     'GATCGTAGCTAGTTTTTTTA');
ANALYZE tutorial_mmsearch;
\end{lstlisting}
\end{minipage}

If rows were loaded with maintenance off, first run
\texttt{bioseq.update\_index} with the required recipe as described in
Section~\ref{sec:appendix-indexes}. \texttt{ANALYZE} refreshes PostgreSQL
statistics for row counts and ordinary column distributions. It does not
build minimizer or symbol-count indexes.

\subsection{Keep each alignment as a separate row}

\par\noindent\begin{minipage}{\linewidth}
\begin{lstlisting}[style=bioseq-appendix-sql]
SELECT pk, organism,
       COUNT(*) OVER (PARTITION BY pk) AS matches_for_row,
       bioseq.query_start(pk) AS query_start,
       bioseq.query_end(pk) AS query_end,
       bioseq.target_start(pk) AS target_start,
       bioseq.target_end(pk) AS target_end,
       bioseq.cigar(pk) AS cigar
FROM tutorial_mmsearch
WHERE organism = 'Homo sapiens'
  AND collected_on >= DATE '2025-01-01'
  AND public.mmsearch(
      seq0,
      'CTAGTCACTGTAGCCTGATCGTAGCTAGT'::bioseq.seq('dna4', 64),
      public.aligner_options(best_n => 3)
  )
ORDER BY pk, target_start, target_end;
\end{lstlisting}
\end{minipage}

\begin{itemize}
  \item The metadata predicates select eligible targets. The window count
    counts alignment rows for each eligible target after search filtering.
  \item Each coordinate helper reads the current alignment through the
    search's table alias. The key is an unambiguous carrier because this
    example has one search over one sequence column.
  \item Both intervals are zero-based and half-open. Target 5's literal
    query copies begin at positions 4, 38, and 84; whether all are retained
    depends on the selected algorithms and filtering options.
  \item CIGAR records runs of operations: \texttt{M} consumes symbols on
    both sides, \texttt{I} consumes query symbols, and \texttt{D} consumes
    target symbols. \texttt{29M} describes a 29-symbol aligned run;
    \texttt{M} alone does not distinguish a match from a substitution.
\end{itemize}

Do not collapse alignments by primary key when their separate positions
matter. When only matching target keys are needed, request that different
result explicitly:

\par\noindent\begin{minipage}{\linewidth}
\begin{lstlisting}[style=bioseq-appendix-sql]
SELECT DISTINCT pk
FROM tutorial_mmsearch
WHERE organism = 'Homo sapiens'
  AND public.mmsearch(
      seq0,
      'CTAGTCACTGTAGCCTGATCGTAGCTAGT'::bioseq.seq('dna4', 64)
  )
ORDER BY pk;
\end{lstlisting}
\end{minipage}

\subsection{Choose the result fields}

Scores, edits, and seed counts describe different parts of an alignment.
They are not interchangeable percentages.

\par\noindent\begin{minipage}{\linewidth}
\begin{lstlisting}[style=bioseq-appendix-sql]
SELECT pk,
       bioseq.score(pk) AS score,
       bioseq.chain_score(pk) AS chain_score,
       bioseq.edit(pk) AS edits,
       bioseq.divergence(pk) AS divergence,
       bioseq.seed_count(pk) AS seed_count,
       bioseq.mapq(pk) AS mapq,
       bioseq.is_primary(pk) AS is_primary,
       bioseq.target_start(pk) AS target_start
FROM tutorial_mmsearch
WHERE public.mmsearch(
    seq0,
    'CTAGTCACTGTAGCCTGATCGTAGCTAGT'::bioseq.seq('dna4', 64),
    public.aligner_options(min_chain_score => 30, best_n => 2)
)
ORDER BY score DESC, pk, target_start;
\end{lstlisting}
\end{minipage}

\begin{itemize}
  \item \texttt{score} is an alias for \texttt{dp\_score}, the base-alignment
    score. \texttt{chain\_score} is the earlier seed-chain score.
  \item \texttt{edit} reports the alignment's edit count.
    \texttt{seed\_count} counts anchors retained in its seed chain.
  \item \texttt{divergence} is an estimate computed from seeds; it is not
    the edit count divided by the aligned span. \texttt{mapq} is the
    reported mapping quality, not a measured percentage identity.
  \item \texttt{is\_primary} reports the region classification made during
    search. Several separated primary regions can survive for one target.
\end{itemize}

Keep the effective options with saved results. In particular, scores from
different scoring schemes need not share a useful scale. Request only fields
the client uses: \texttt{cigar}, \texttt{has\_cigar}, and \texttt{edit} request
traceback-related work, while divergence and mapping quality have their own
computation flags. A key-only query can avoid work required only for these
projections.

\subsection{Bind parameters and combine metadata filters}

This prepared statement binds an organism, a date, and query text. The cast
checks query symbols against the target alphabet. It does not recode a
different bound alphabet automatically.

\par\noindent\begin{minipage}{\linewidth}
\begin{lstlisting}[style=bioseq-appendix-sql]
PREPARE find_tutorial_matches(text, date, text) AS
SELECT pk,
       bioseq.target_start(pk) AS target_start,
       bioseq.target_end(pk) AS target_end,
       bioseq.cigar(pk) AS cigar
FROM tutorial_mmsearch
WHERE organism = $1
  AND collected_on >= $2
  AND public.mmsearch(seq0, $3::bioseq.seq('dna4', 64))
ORDER BY pk, target_start, target_end;

EXECUTE find_tutorial_matches(
    'Homo sapiens', DATE '2025-01-01',
    'CTAGTCACTGTAGCCTGATCGTAGCTAGT'
);
DEALLOCATE find_tutorial_matches;
\end{lstlisting}
\end{minipage}

The query sequence may be a runtime parameter. The options object has a
different restriction: it must be reducible to a constant when planned.
Section~\ref{sec:appendix-search-options} gives the complete object API and
its precedence over GUC defaults. In the remaining listings, \texttt{\$1}
denotes query text bound through a prepared statement or the PostgreSQL
extended query protocol; these are application statement templates.

\subsection{Separate search limits from SQL limits}

With a positive \texttt{pri\_ratio}, \texttt{best\_n} limits retained
secondary regions across the complete region list passed to the
secondary-filtering step. Primary regions do not count towards it, and the
limit is not applied separately to each primary region.
Setting \texttt{pri\_ratio} to zero bypasses
this secondary-selection step, including that count check. Therefore
\texttt{best\_n => 3} is not a promise of three output rows.

Use SQL \texttt{LIMIT} for a limit on the final alignment rows, after an
explicit ordering:

\par\noindent\begin{minipage}{\linewidth}
\begin{lstlisting}[style=bioseq-appendix-sql]
SELECT pk,
       bioseq.target_start(pk) AS target_start,
       bioseq.target_end(pk) AS target_end,
       bioseq.cigar(pk) AS cigar
FROM tutorial_mmsearch
WHERE public.mmsearch(
    seq0, $1::bioseq.seq('dna4', 64),
    public.aligner_options(best_n => 3)
)
ORDER BY pk, target_start, target_end
LIMIT 100;
\end{lstlisting}
\end{minipage}

Ordering by key and both endpoints gives predictable ordering for distinct
coordinate pairs. If other alignment fields can differ at identical
coordinates, include those fields as additional ordering keys before
implementing pagination.

\subsection{Change thresholds and request an edit prefilter}

The object groups parameters for minimizers, chaining, base-alignment
scoring, secondary selection, and execution choices. Change one group at a
time. A higher minimum chain score removes weaker seed chains before base
alignment. Changing minimizer parameters can also make a stored index
incompatible: an allowing read policy can recompute, whereas strict persisted
mode requires a compatible stored generation.

The recognized same-level edit predicate has this exact form:

\par\noindent\begin{minipage}{\linewidth}
\begin{lstlisting}[style=bioseq-appendix-sql]
SELECT pk, bioseq.edit(seq0) AS edits
FROM tutorial_mmsearch
WHERE public.mmsearch(seq0, $1::bioseq.seq('dna4', 64))
  AND bioseq.edit(seq0) < 5
ORDER BY pk;
\end{lstlisting}
\end{minipage}

\begin{itemize}
  \item Use the sequence column as the edit predicate's argument. A key
    carrier accepted by an output helper does not satisfy this predicate's
    stricter contract.
  \item Keep exactly one top-level comparison \texttt{edit(sequence) <}
    an integer constant. The exclusive bound 5 admits edit counts below 5.
  \item The prefilter can reject candidates before full alignment, but
    surviving results still undergo alignment and the final edit test.
    Inspect execution counters to establish which prefilter actually ran.
  \item A transformed target, such as a slice, cannot use this same-level
    edit predicate. Project edits inside the search query and filter that
    ordinary result column in a separate materialized query level.
\end{itemize}

\subsection{Calculate spans and restore slice coordinates}

First project the match helpers as ordinary columns. Then perform arithmetic
in an outer query. \texttt{MATERIALIZED} preserves that boundary.

\par\noindent\begin{minipage}{\linewidth}
\begin{lstlisting}[style=bioseq-appendix-sql]
WITH matches AS MATERIALIZED (
    SELECT pk,
           bioseq.query_start(pk) AS qs,
           bioseq.query_end(pk) AS qe,
           bioseq.target_start(pk) AS ts,
           bioseq.target_end(pk) AS te
    FROM tutorial_mmsearch
    WHERE public.mmsearch(seq0, $1::bioseq.seq('dna4', 64))
)
SELECT pk, qs, qe, qe - qs AS query_span,
       ts, te, te - ts AS target_span
FROM matches
ORDER BY pk, ts, te;
\end{lstlisting}
\end{minipage}

Do not add one: interval \([10,15)\) contains five positions. For a sliced
target, helper coordinates refer to the slice. Use the same slice expression
in the helpers and the search, then add the slice's resolved origin.

\par\noindent\begin{minipage}{\linewidth}
\begin{lstlisting}[style=bioseq-appendix-sql]
WITH relative_matches AS MATERIALIZED (
    SELECT pk,
           bioseq.target_start(seq0[4:100]) AS ts,
           bioseq.target_end(seq0[4:100]) AS te
    FROM tutorial_mmsearch
    WHERE public.mmsearch(
        seq0[4:100], $1::bioseq.seq('dna4', 64)
    )
)
SELECT pk, 4 + ts AS stored_start, 4 + te AS stored_end
FROM relative_matches
ORDER BY pk, stored_start, stored_end;
\end{lstlisting}
\end{minipage}

The fixed lower bound 4 is inside every non-NULL target in this fixture. A
negative or row-dependent bound must first be resolved and clamped against
that row's length; adding the literal negative bound would be wrong.

\subsection{Inspect a search plan and compare routes}

Explain the final statement with its real projection list. Begin with
\texttt{EXPLAIN (VERBOSE, COSTS OFF)}. When execution is acceptable, add
\texttt{ANALYZE}; JSON output is convenient for storing selected plan fields.

\par\noindent\begin{minipage}{\linewidth}
\begin{lstlisting}[style=bioseq-appendix-sql]
EXPLAIN (ANALYZE, VERBOSE, COSTS OFF, FORMAT JSON)
SELECT pk
FROM tutorial_mmsearch
WHERE organism = 'Homo sapiens'
  AND public.mmsearch(
      seq0,
      'CTAGTCACTGTAGCCTGATCGTAGCTAGT'::bioseq.seq('dna4', 64)
  );
\end{lstlisting}
\end{minipage}

Read the selected custom node, projected columns, index-read requirement,
planned index/recompute mix, pipeline, and Highway dispatch. With
\texttt{ANALYZE}, inspect the observed target-index source and row counts,
backend, GPU dispatches, and executed CPU architecture when a target ran.
Some nodes print runtime fields before execution, so their presence alone
does not establish that any work ran. Compare estimates with alignment rows:
one target can yield several alignments. Refresh ordinary column statistics
when metadata estimates are poor; rebuilding minimizers does not fix stale
scalar statistics.
Section~\ref{sec:appendix-explain} provides the full plan-field reference.

For a controlled index-read comparison under an allowing read policy, repeat
the same search with stored-index reads enabled and disabled:

\par\noindent\begin{minipage}{\linewidth}
\begin{lstlisting}[style=bioseq-appendix-sql]
SET bioseq.index.sequence_read.enabled = on;
-- Run and explain the search; save its ordered alignments.
SET bioseq.index.sequence_read.enabled = off;
-- Repeat the same search, using recomputed target data.
RESET bioseq.index.sequence_read.enabled;
\end{lstlisting}
\end{minipage}

The switch permits or forbids reading stored sequence-index chunks. Strict
per-call requirements can reject a conflicting configuration. Compare actual
index source, persisted rows, missing persisted rows, and recomputed rows.
Check matching coordinates and result fields before interpreting timings.

Similarly, repeat the statement with vectorized SQL execution disabled and
enabled:

\par\noindent\begin{minipage}{\linewidth}
\begin{lstlisting}[style=bioseq-appendix-sql]
SET bioseq.vectorization.enabled = off;
-- Run and explain the application statement.
SET bioseq.vectorization.enabled = on;
-- Repeat it and compare the selected pipeline and results.
RESET bioseq.vectorization.enabled;
\end{lstlisting}
\end{minipage}

This Boolean permits vectorized SQL predicate and aggregate execution where
the plan supports it. Some projection and lazy-reading paths remain scalar.
It does not choose one CPU instruction set or independently select every
alignment algorithm. Highway chooses among compiled implementations
supported by the CPU. If a comparison
also changes the minimizer scheme or uses bounded SIMD chaining, identical
alignment rows are not a general guarantee. Record the effective algorithms
and compare results before measuring speed. Vulkan availability likewise does
not prove a GPU stage ran; use actual backend and dispatch counters.
Section~\ref{sec:appendix-highway} explains CPU dispatch separately.

\subsection{Search ownership, errors, and cleanup}

Two independent searches over the same target column and table alias in one
\texttt{WHERE} clause cannot share one set of match helpers. The following
shape is rejected:

\par\noindent\begin{minipage}{\linewidth}
\begin{lstlisting}[style=bioseq-appendix-sql]
-- Unsupported: two independent searches share one target alias.
SELECT pk FROM tutorial_mmsearch
WHERE public.mmsearch(seq0, $1::bioseq.seq('dna4', 64))
  AND public.mmsearch(seq0, $2::bioseq.seq('dna4', 64));
\end{lstlisting}
\end{minipage}

Use separate statements and intersect their keys, or distinct table aliases
in a supported join. Section~\ref{sec:appendix-joins} shows joins and self-joins with
separate search owners. Keep the match helpers at the query level where their
search is defined.

\begin{itemize}
  \item \texttt{22P02}: query text contains a symbol rejected by its bound
    alphabet, such as N in \texttt{dna4}.
  \item \texttt{42804}: target and query have incompatible bound alphabets.
  \item \texttt{0A000}: an unsupported search or helper shape, including
    stacked searches or an unsupported edit predicate.
  \item \texttt{42883}: a function or operator could not be resolved. Check
    signatures, installation schema, and search path.
  \item Zero result rows with valid input mean no retained alignment. This
    is different from any of these errors.
\end{itemize}

If a search has no custom owner, inspect its access method, operand types,
predicate shape, and route settings before adjusting thresholds. After
maintenance-off writes, rebuild the selected index before requiring persisted
coverage. A fresh connection, warm-up query, and stable result ordering are
useful when comparing routes, but a tiny first query is not a throughput
measurement.

\par\noindent\begin{minipage}{\linewidth}
\begin{lstlisting}[style=bioseq-appendix-sql]
DROP TABLE tutorial_mmsearch;
RESET bioseq.index.insert_maintenance.enabled;
\end{lstlisting}
\end{minipage}
